\backhead{FURTHER EXERCISES}{khred}
\begin{exlist}
\item A log gives a storage organ dry matter of 7.2 Mg DM\,ha$^{-1}$. Find the yield at 14 per cent moisture.\exans{8.37 Mg\,ha$^{-1}$}
\item Close the balance: $P=540$, $I=120$, $T=290$, $E_s=150$, $E_o=45$, $D=160$, $Q=5$ mm, $\Delta S=-6$ mm.\exans{residual 16 mm: a term is missing}
\item A daily log of 20 years has how many rows? An hourly log?\exans{7305; 175\,320}
\item Convert 0.85 g\,cm$^{-3}$ to kg\,m$^{-3}$ and 12 kg N\,ha$^{-1}$ to g N\,m$^{-2}$.\exans{850 kg\,m$^{-3}$; 1.2 g\,m$^{-2}$}
\item The transpiration efficiency is 2.1 kg DM\,m$^{-3}$ and the transpiration 280 mm. Find the dry matter.\exans{5.88 Mg\,ha$^{-1}$}
\end{exlist}
\backhead{ESSAY QUESTIONS}{khred}
{\small\itshape Answer each question in about one page, with equations, a sketch and at least one number from this chapter.}\par\vspace{3pt}
\begin{exlist}
\item Describe the structure of the Daisy output and how it is read into SWAFL.
\item Discuss the checks that must be made before a log is used.
\item Explain why daily figures are drawn only from some experiments.
\item Discuss the conversion of dry matter into yield and water productivity.
\item Explain how a failed run is recognised and what can be learnt from it.
\item Describe how you would present the output of a Daisy experiment in a report.
\end{exlist}
